% ================================
% === Reviewer 1 ===
% ================================
\reviewerblock{Reviewer~1}{%
We thank the reviewer for a very constructive review and for recognizing the importance of the problem and the empirical interest of FDIL.}

\cresp{The main concerns are the fairness and interpretability of the comparisons. Please provide controlled baselines under the same backbone and thresholding protocol.}
{This is an important fairness concern. We change both the protocol description and the result presentation. In Experimental Setup, we now define the controlled setting explicitly: FDIL, the CLIP baseline, HLEG, LabCR, and IntCLIP use the same frozen CLIP ViT-L/14 features, data split, and validation-only threshold search. In Table~\ref{tab:main_results}, we add lines with ``(CLIP ViT-L/14)'' from the methods' original-pipeline results, which are retained only as a complementary robustness check. Under the matched protocol, FDIL leads the three representative external methods by \res{$+3.4$ to $+4.7$ Avg F1}. Finally, we add an Interpretability Analysis subsection and a per-class attribution table: the text identifies the largest gains in abstract, disagreement-prone intents and explicitly reports the only two regressions among the 28 classes.

\revloc{Experimental Setup and Comparison with state-of-the-art methods in Sec.~\ref{sec:experiments}; the matched-feature rows and original-pipeline rows in Table~\ref{tab:main_results}; Interpretability Analysis and Table~\ref{tab:per_class_attr}; the claim boundary in Sec.~\ref{sec:discussion} (Limitations and Future Directions)}
\begin{table}[h]
\centering
\caption{Per-class F1 attribution of FDIL versus the CLIP baseline under the
matched class-wise protocol. Results use a single representative seed. The
table reports the five largest improvements and the only two regressions among
all 28 classes.}
\label{tab:per_class_attr}
\begin{tabular}{lcccc}
\toprule
Intent class & Subset & Baseline & FDIL & $\Delta$ \\
\midrule
WorkILike & Hard & 14.6 & 43.0 & $+28.5$ \\
PassionAboutSomething & Hard & 11.3 & 28.6 & $+17.3$ \\
SocialLifeFriendship & Hard & 33.9 & 51.2 & $+17.2$ \\
ManageableMakePlan & Medium & 17.5 & 32.0 & $+14.5$ \\
BeatCompete & Medium & 54.6 & 68.1 & $+13.5$ \\
\midrule
FineDesignLearnArt-Arch & Easy & 84.0 & 83.1 & $-1.0$ \\
CreativeUnique & Medium & 41.9 & 36.6 & $-5.3$ \\
\bottomrule
\end{tabular}
\end{table}
}
\showMainResults

\cresp{The strong effect of class-wise thresholding.}
{To separate calibration gain from method gain, we add a dedicated ``Calibration and Threshold-free Evaluation'' subsection. Its first table is a four-cell, three-seed comparison crossing method (CLIP baseline versus FDIL) with threshold type (global versus class-wise). This shows that FDIL improves Avg F1 under both class-wise (\res{$+4.08$}) and global (\res{$+2.31$}) thresholding, while class-wise calibration benefits both methods. We then add a second table with four threshold-free metrics and a 5,000-resample paired bootstrap test, plus ROC and precision--recall curves computed from the same scores in Fig.~\ref{fig:roc_pr}. FDIL improves over the CLIP baseline on every threshold-free metric; its mean mAP gain is \res{$+4.43$}, and the representative-seed difference is \res{$+4.80$} with 95\% CI \res{$[3.58,6.10]$} and $p<0.001$.
\revloc{Calibration and Threshold-free Evaluation in Sec.~\ref{sec:experiments}; the method$\times$threshold cells in Table~\ref{tab:thresholding_test}; the mAP, ROC-AUC, PR-AUC, confidence-interval, and significance columns in Table~\ref{tab:threshold_free}; ROC/PR curves in Fig.~\ref{fig:roc_pr}}
}
\showThresholdingTest
\showThresholdFree
\showRocPr

\cresp{Report variance over runs.}
{Following this suggestion, we rerun the UTD-only and full FDIL settings with three seeds. We replace the former single-run entries in the main comparison with mean$\pm$std and use the same three runs for every cell of the thresholding decomposition, so the reported variance reflects training randomness rather than different evaluation protocols.
\revloc{Experimental Setup in Sec.~\ref{sec:experiments} (seed protocol); the UTD and FDIL rows in Table~\ref{tab:main_results}; all four cells in Table~\ref{tab:thresholding_test}}
}
\showMainResults
\showThresholdingTest

\cresp{Clarify validation-only model selection.}
{We make the selection boundary explicit and audit the implementation. In the Implementation Details paragraph, we now state the sequence of operations: early stopping and hyperparameter selection use validation metrics only; global or class-wise thresholds are fitted on the validation split; those thresholds are then frozen; and the test split is evaluated once for final reporting. We also add an explicit sentence that test outputs never enter a model- or threshold-selection branch.
\revloc{Implementation Details in Sec.~\ref{sec:experiments}, in the paragraph describing early stopping, threshold fitting, and final test evaluation}
}

\cresp{The dependence on offline VLM-generated rationales and semantic priors, and incomplete reproducibility details. Please release or document generated artifacts.}
{We now state this dependence and its cost explicitly at three locations. First, the SLR-C and UTD subsections identify which inputs are generated offline and clarify that the UTD text teacher is discarded after training, and the full FDIL model retains the lightweight semantic-prior reranking stage. Second, Implementation Details now records the reproducible pipeline order, which is feature extraction, prior construction, teacher encoding, optimization, validation-only threshold search, and agreement gating. Third, we replace ``available on request'' in the availability statement with an itemized release commitment covering generated priors and rationales, cached features/logits, exact prompts, fitted thresholds, seeds, checksums, and training/evaluation scripts upon acceptance.
\revloc{Semantic Prior Construction and Text-derived Teacher in Sec.~\ref{sec:method}; the pipeline paragraph in Implementation Details of Sec.~\ref{sec:experiments}; the final Data and artifact availability statement}
}

\cresp{Tone down claims about functional decoupling and SOTA performance unless the ablations isolate these factors more cleanly.}
{We agree that these claims require evidence that separates architectural organization from backbone and calibration effects. We therefore revise the manuscript at the following specific locations.

First, in the Abstract, we replace the broad performance statement in the final sentence with protocol-qualified effect sizes: FDIL improves over representative recent methods by at least $+3.4$ Avg F1, $+5.7$ Hard F1, and $+3.4$ threshold-free mAP under matched frozen CLIP ViT-L/14 features and validation-only thresholding.

Second, in ``Comparison with state-of-the-art methods'' of Sec.~\ref{sec:experiments}, we add a fairness caveat immediately after Table~\ref{tab:main_results}. The revised paragraph explicitly states that results quoted from methods using their original, weaker backbones are not evidence of FDIL's algorithmic superiority. We therefore add separate HLEG, LabCR, and IntCLIP rows marked ``CLIP ViT-L/14'' to Table~\ref{tab:main_results}; these rows use the same frozen features, split, and validation-only class-wise thresholds as FDIL. Under this controlled comparison, FDIL retains a \res{$+3.4$ to $+4.7$ Avg F1} advantage.

Third, in the Ablation Studies, we add the new subsection ``Unified versus Functionally Decoupled Ambiguity Handling'' and Table~\ref{tab:unified_decoupled}. The most direct unified control averages the semantic-prior and rationale-teacher probabilities into one distillation target and trains a single student under the same features and selection rules. It reaches \res{$55.66$ Avg F1}, compared with \res{$57.49$} for decoupled FDIL. We also add ``Negative Controls on the Text Teacher'' and Table~\ref{tab:negative_controls}, where shuffled and generic rationales test whether the gain comes merely from adding external text, and gate-removal/flattening variants test the role of agreement-aware supervision.

Fourth, in ``Calibration and Threshold-free Evaluation'', we add Table~\ref{tab:thresholding_test} to cross the CLIP baseline and FDIL with global versus class-wise thresholds over three seeds. We added Table~\ref{tab:threshold_free} and Fig.~\ref{fig:roc_pr} to report mAP, ROC-AUC, PR-AUC, micro-averaged ROC/PR curves, and a 5,000-resample paired bootstrap test. These additions show that the FDIL gain remains when class-wise thresholding is removed and when evaluation is made fully threshold-free.

Finally, in the Discussion, we add a bounded interpretation of the evidence: functional decoupling provides a measurable aggregate benefit and keeps semantic reranking and supervision stabilization separately interpretable and deployable, but the unified controls remain competitive on some metrics. Accordingly, we no longer attribute improvements observed across incompatible original pipelines to the proposed organization itself.
\revloc{the final sentence of the Abstract; Comparison with state-of-the-art methods and the matched-feature rows in Table~\ref{tab:main_results}; Unified versus Functionally Decoupled Ambiguity Handling and Table~\ref{tab:unified_decoupled}; Negative Controls on the Text Teacher and Table~\ref{tab:negative_controls}; Calibration and Threshold-free Evaluation, Tables~\ref{tab:thresholding_test} and~\ref{tab:threshold_free}, and Fig.~\ref{fig:roc_pr}; the interpretation in Sec.~\ref{sec:discussion}}
}
\showMainResults
\showUnifiedDecoupled
\showNegativeControls
\showThresholdingTest
\showThresholdFree
\showRocPr
